
به ازای هر راس میانگین درجه برابر با $d = o(\log n)$ هست. 
در نتیجه $p = o(\frac{\log n}{n})$. درجه هر راس از توزیع زیر پیروی میکند :
\begin{latin}
\begin{equation*}
    P(\text{deg}(v) = k) = {n-1\choose k} p^k(1-p)^{n-1-k}
\end{equation*}
\end{latin}
به ازای $n$ های به قدر کافی بزرگ و $p$ های به قدر کافی کوچک میتوانیم توزیع دو جمله ای بالا را با یک توزیع پوآسون با پارامتر $d = o(np)$ تقریب بزنیم. پس :
\begin{latin}
\begin{equation*}
    p(\text{deg}(v) = k) \sim \frac{d^k\exp(-d)}{k!}
\end{equation*}

\begin{flalign*}
    \Rightarrow \mathbb{P}\big\lceil \exists v,\text{deg}(v) = 10d\big\rceil &= 1- \mathbb{P}\big\lceil \forall V, \text{deg}(v) \neq 10d \big\rceil &\\
    &= 1 - \mathbb{P}\big\lceil \text{deg}(v) \neq 10d \big\rceil^n\\
    &= 1 - \big\lceil 1-\mathbb{P}\big\lceil \text{deg}(v) = 10d \big\rceil \big\rceil^n\\
    &= 1 - (1 - \frac{d^{10d} \exp(-d)}{10d!})^n\\
    &\approx 1 - (1 - \frac{1}{\sqrt{2\pi 10d }}(\frac{e}{10d})^{10d}\cdot d^{10d}\exp(-d))^n \quad(\text{Striling})\\
    &= 1 - (1 - \frac{1}{\sqrt{2\pi 10d }}\frac{e}{10d}^{10d}\exp(-d) )^n
\end{flalign*}
\begin{equation*}
    \textcolor{red}{**} d = o(\log n) \Rightarrow d\leq C\log n \Rightarrow \frac{1}{\sqrt{2\pi 10d }}(\frac{e}{10d})^{10d}\exp(-d)\leq (\frac{e}{10})^{C\log(n)} \times n^{-c'} \times C" \leq n^{c-c'} \times k \leq n^{-M}
\end{equation*}
\begin{flalign*}
    &\leq 1-(1-n^{-M})^n
\end{flalign*}

\end{latin}

\begin{tcolorbox}[colback=yellow!10!white,colframe=red!75!black]
    \begin{equation*}
        n,m > 0 \Rightarrow 0 < n^{-M} < 1 \Rightarrow (1- n^{-M})^n \xrightarrow{n\rightarrow \infty} 0
    \end{equation*}
    \begin{latin}
    \begin{equation*}
        \Rightarrow \text{W.H.P} : \mathbb{P}\big\lceil \exists v, \text{deg}(v) = 10d \big\rceil
    \end{equation*}
    \end{latin}
\end{tcolorbox}
\subsection{بخش دوم}
در این مسئله یک گراف تصادفی
$G \sim G(n,p)$
داریم که درجه‌ی مورد انتظار هر رأس در آن
\[
d = (n-1)p = O(1)
\]
است. بنابراین عدد ثابتی $M>0$ وجود دارد که برای $n$های بزرگ
\[
d \le M
\]
برقرار است. درجه‌ی هر رأس $i$ را با $d_i$ نشان می‌دهیم. می‌دانیم
\[
d_i \sim \operatorname{Bin}(n-1,p)
\]
و امیدریاضی آن برابر
\[
\mathbb{E}[d_i] = (n-1)p = d
\]
است.

برای متغیر دوجمله‌ای از نابرابری چرنوف استفاده می‌کنیم. اگر
$X \sim \operatorname{Bin}(N,p)$
با میانگین $\mu = Np$ باشد و $k \ge \mu$، آنگاه
\[
\Pr(X \ge k) \le e^{-\mu} \left( \frac{e\mu}{k} \right)^k .
\]

این کران را روی $d_i$ اعمال می‌کنیم. آستانه را به صورت
\[
k = C \cdot \frac{\log n}{\log\log n}
\]
در نظر می‌گیریم که در آن $C>0$ یک ثابت است که بعداً مقدارش را انتخاب می‌کنیم. در نتیجه
\begin{align*}
\Pr\!\left(d_i \ge C \frac{\log n}{\log\log n}\right)
&\le e^{-d}
\left(
\frac{e d}{C \frac{\log n}{\log\log n}}
\right)^{C \frac{\log n}{\log\log n}} \\
&= e^{-d}
\left(
\frac{e d \log\log n}{C \log n}
\right)^{C \frac{\log n}{\log\log n}}.
\end{align*}
از آن‌جا که $d \le M$، می‌توانیم $d$ را با $M$ جایگزین کنیم و به کرانی ساده‌تر برسیم:
\[
\Pr\!\left(d_i \ge C \frac{\log n}{\log\log n}\right)
\le
e^{-d}
\left(
\frac{e M \log\log n}{C \log n}
\right)^{C \frac{\log n}{\log\log n}}.
\]

حال احتمال این‌که «حداقل یک رأس» درجه‌ای بزرگ‌تر از این آستانه داشته باشد را با استفاده از کران اتحادی \lr{(union bound)} می‌بندیم:
\begin{align*}
\Pr\Big(\exists\, i:\ d_i \ge C \frac{\log n}{\log\log n}\Big)
&\le
\sum_{i=1}^n
\Pr\!\left(d_i \ge C \frac{\log n}{\log\log n}\right) \\
&\le
n \, e^{-d}
\left(
\frac{e M \log\log n}{C \log n}
\right)^{C \frac{\log n}{\log\log n}}.
\end{align*}
اگر این عبارت را $T_n$ بنامیم، داریم
\[
T_n =
n \, e^{-d}
\left(
\frac{e M \log\log n}{C \log n}
\right)^{C \frac{\log n}{\log\log n}}.
\]
کافی است نشان بدهیم
\[
T_n \xrightarrow[n\to\infty]{} 0,
\]
چون در این صورت با احتمال میل‌کننده به یک، هیچ رأسی درجه‌ی بزرگ‌تر از
$C \frac{\log n}{\log\log n}$
نخواهد داشت و در نتیجه حداکثر درجه‌ی گراف از همین مرتبه است.

برای بررسی حد، لگاریتم می‌گیریم. اگر
$\log T_n \to -\infty$،
آنگاه $T_n \to 0$. بنابراین
\[
\log T_n
=
\log n - d
+
C \frac{\log n}{\log\log n}
\log\!\left(
\frac{e M \log\log n}{C \log n}
\right).
\]
برای ساده‌سازی، ثابت را
$C = M$
انتخاب می‌کنیم. آن‌گاه
\begin{align*}
\log T_n
&=
\log n - d
+
M \frac{\log n}{\log\log n}
\log\!\left(
\frac{e M \log\log n}{M \log n}
\right) \\
&=
\log n - d
+
M \frac{\log n}{\log\log n}
\log\!\left(
\frac{e \log\log n}{\log n}
\right).
\end{align*}
لگاریتم داخل پرانتز را باز می‌کنیم:
\[
\log\!\left(
\frac{e \log\log n}{\log n}
\right)
= \log e + \log(\log\log n) - \log(\log n)
= 1 + \log(\log\log n) - \log(\log n).
\]
در نتیجه
\[
\log T_n
=
\log n - d
+
M \frac{\log n}{\log\log n}
\big(1 + \log(\log\log n) - \log(\log n)\big).
\]

حال رفتار جمله‌ی داخل پرانتز را وقتی $n$ بزرگ می‌شود بررسی می‌کنیم. چون
$\log(\log n)$
خیلی سریع‌تر از
$\log(\log\log n)$
رشد می‌کند، داریم
\[
1 + \log(\log\log n) - \log(\log n) \sim - \log(\log n)
\]
و بنابراین این جمله برای $n$های بزرگ منفی و با قدر مطلق بزرگ می‌شود. این مقدار منفی در
$M \frac{\log n}{\log\log n}$
ضرب می‌شود و تقریباً رفتاری شبیه
\[
- M \log n
\]
دارد. پس می‌توان نوشت
\[
\log T_n
\approx
\log n - d - M \log n
= (1 - M)\log n - d.
\]
از آن‌جا که $M$ را می‌توانیم بزرگ‌تر از $1$ انتخاب کنیم، ضریب
$1 - M$
منفی است و در نتیجه
$(1 - M)\log n$
وقتی $n \to \infty$ به $-\infty$ میل می‌کند و عدد ثابت $-d$ در مقابل آن ناچیز است. بنابراین
\[
\log T_n \xrightarrow[n\to\infty]{} -\infty
\quad\Rightarrow\quad
T_n \xrightarrow[n\to\infty]{} 0.
\]

به این ترتیب احتمال «وجود یک رأس با درجه‌ی بزرگ‌تر از
$C \frac{\log n}{\log\log n}$
» به صفر میل می‌کند؛ یعنی
\[
\Pr\Big(\forall\, i:\ d_i \le C \frac{\log n}{\log\log n}\Big)
\xrightarrow[n\to\infty]{} 1.
\]
پس در گراف تصادفی $G(n,p)$ با درجه‌ی مورد انتظار ثابت، حداکثر درجه‌ی گراف با احتمال بالا
\[
O\!\left(\frac{\log n}{\log\log n}\right)
\]
است.
